\section{Measured comparison}
\label{sec:results}

All figures are on one Apple~M4 (10-core, $16$\,GB), non-zero-knowledge,
provers multithreaded, at the matched commitment rate $1/4$. \zinc{} instances
are indexed by $\nu=\log_2(\text{trace rows})$ with
$\mathrm{compressions}(\nu)=(2^\nu-4)/68$; \binius{} by SHA-256 block count,
matched to \zinc{}'s compression count. The detailed sweeps --- prover,
verifier, proof, and the per-phase decompositions --- live in \Cref{sec:impl}
(\Cref{tab:scaling,tab:perphase,tab:perphase-b64}); this section reads off the
comparison they imply. Absolute times are thermally variable
($\sim$$10$--$20\%$ run-to-run), so we lean on ratios and per-phase
\emph{shapes}, which are stable. The Zinc\texttt{+} figures here are CPU-only
(no GPU commit), so the prover comparison is core-for-core; \binius{}'s
verifier carries a parallelised public-input sweep (a constant-factor patch on
an $O(N)$ algorithm).

\subsection{Headline}
\label{sec:res-headline}

\begin{table}[t]
\centering
\renewcommand{\arraystretch}{1.2}
\begin{tabular}{@{}llrrl@{}}
\toprule
& & \binius{} & \zinc{} & winner \\
\midrule
\multirow{3}{*}{$\sim$$1000$ comp}
  & Prove   & $84$\,ms    & $75$\,ms    & \zinc{} $1.1\times$ \\
  & Verify  & $14$\,ms    & $12$\,ms    & \zinc{} $1.2\times$ \\
  & Proof   & $249$\,KiB  & $1.12$\,MiB & \binius{} $4.6\times$ \\
\midrule
\multirow{3}{*}{$\sim$$8000$ comp}
  & Prove   & $1.22$\,s   & $0.58$\,s   & \zinc{} $2.1\times$ \\
  & Verify  & $493$\,ms   & $45$\,ms    & \zinc{} $11\times$ \\
  & Proof   & $299$\,KiB  & $\sim$$2.0$\,MiB & \binius{} $\sim$$7\times$ \\
\bottomrule
\end{tabular}
\caption{Matched rate $1/4$, CPU, at two scales (full curves:
\Cref{tab:scaling}). At moderate $N$ the prover is close ---
\zinc{} a touch ahead at the matched rate (the GPU commit tracks the CPU
prover without a net latency win at this scale, \Cref{tab:scaling-gpu});
at scale the gaps open. The prover divergence is
\binius{}'s super-linearity (\Cref{tab:perphase-b64}) against \zinc{}'s linear
scaling (\Cref{tab:perphase}); the verifier divergence is $O(N)$ vs $\sqrt N$;
the proof runs the other way (recursive FRI vs the flat bit-sliced opening).
Representative single runs.}
\label{tab:headline}
\end{table}

\Cref{tab:headline} is the comparison in two regimes. At a matched
$\sim$$1000$ compressions all three numbers are close-to-modest: \zinc{}'s
prover is a touch ahead at rate $1/4$ (its GPU commit tracks the CPU prover
without a net latency win at this scale, \Cref{tab:scaling-gpu}), the
verifiers are within $\sim$$1.2\times$, and \binius{}'s proof is several
times smaller. The interesting content is how they \emph{scale}. \zinc{}'s
prover is linear and \binius{}'s is super-linear, so the prover gap opens to
$\sim$$2\times$ by $8000$ compressions and then becomes unbounded: \binius{}'s
$O(N)$ circuit \emph{build} exceeds one minute beyond $\sim$$8000$
compressions, where \zinc{} still proves $15{,}420$ in $\approx 1.2$\,s. The
verifier and proof move in opposite directions with $N$, for the structural
reasons of \Cref{sec:axes}.

\subsection{Per-phase: a net of offsetting differences}
\label{sec:res-phases}

The prover comparison is a \emph{net} of large, offsetting per-phase
differences, not a uniform gap. \Cref{tab:perphase,tab:perphase-b64} give the
two decompositions; reading them against each other at the matched instance:

\begin{itemize}
  \item \textbf{Nonlinearity discharge.} \zinc{}'s Hadamard discharge
  ($\sim$$19$\,ms at $\nu{=}16$) is cheaper than \binius{}'s BitAnd $+$ Shift
  reductions ($\sim$$85$\,ms at $\sim$$1000$ comp) --- and it scales better.
  \binius{}'s \emph{Shift reduction} dominates at small $N$ (the entry-wise
  shift handling that \zinc{} gets for free via the $\psi$ read-offs,
  \Cref{sec:shifts}), but at scale the \emph{BitAnd reduction} overtakes it and
  drives the super-linearity (its witness build alone is a quarter of the
  \binius{} prover). \zinc{}'s discharge stays a flat $\sim$$22\%$.
  \item \textbf{Commit.} \binius{}'s FRI commit is cheaper at moderate $N$;
  \zinc{}'s rate-$1/4$ codeword $+$ bit-sliced Merkle cost more (the largest
  \zinc{} phase at scale, $7\%\!\to\!30\%$ of the prover) and is the phase the
  GPU commit most relieves.
  \item \textbf{Open.} \zinc{}'s flat open is \emph{cheaper in time} than
  \binius{}'s FRI open, despite being far larger in \emph{bytes} --- the
  flat-versus-recursive tradeoff is time-versus-size, and \zinc{}'s open scales
  sub-linearly (its $t=987$ openings are a fixed count).
  \item \textbf{Plumbing.} \zinc{} carries the linear PIOP and multipoint
  machinery for which \binius{} has no analogue (its IntMul is a $\le\!1$\,ms
  no-op for SHA).
\end{itemize}

\subsection{Verifier, proof, and preprocessing scale apart}
\label{sec:res-scaling}

Three quantities diverge with $N$, all for the structural reasons of
\Cref{sec:axes}.

\paragraph{Verifier ($O(N)$ vs $\sqrt N$).} \binius{}'s verifier is dominated
by the in-the-clear public-input ``monster'' sweep (\Cref{sec:shifts}): even
with its parallelised verifier it goes $14 \to 493$\,ms from $\sim$$1000$ to
$\sim$$8000$ compressions ($\sim$$35\times$ for $8\times$ work --- super-linear,
as the constraint-index tensor outgrows cache), and is $99\%$ of verify at
scale. \zinc{}'s verifier is $\sqrt N$-bound by its $t$-column proximity test
($12 \to 83$\,ms from $\nu{=}16$ to $\nu{=}20$). The ratio is near-parity at
moderate $N$ and grows past $11\times$ by $8000$ compressions, widening
without bound thereafter.

\paragraph{Proof (recursive vs flat).} The proof moves the other way:
\binius{}'s recursive FRI proof grows slowly ($123\!\to\!299$\,KiB over
$16\!\to\!8192$ compressions), while \zinc{}'s flat bit-sliced opening is
$\sqrt N$ and several megabytes ($0.80\!\to\!2.57$\,MiB at $\nu{=}9\!\to\!20$).
\binius{} wins proof size by $4$--$7\times$, widening with $N$.

\paragraph{Preprocessing ($O(N)$ vs $O(1)$).} The starkest difference is the
circuit description each system must hold (and the verifier must, in effect,
read). \binius{} is a non-uniform word circuit: its constraint system plus
shift/wiring key-collection grow $O(N)$ --- $54.7 / 218 / 873$\,MB at
$256 / 1024 / 4096$ compressions (\Cref{sec:impl-perphase}), the
$\ge\!16{,}384$ build exceeding a minute. \zinc{}'s uniform AIR has $O(1)$
preprocessing: a fixed UAIR specification plus the RAA code's seeds,
kilobyte-scale and independent of the compression count. This
$O(N)$-versus-$O(1)$ gap is the common root of \binius{}'s linear verifier
\emph{and} its circuit-build wall.

\subsection{A correction the measurements forced}
\label{sec:res-commit}

Two earlier figures did not survive a careful re-measurement, and we record the
corrections. First, the \zinc{} verifier at $\nu{=}16$ is $\approx 12$\,ms (the
$t$-column proximity test), not the $\approx 4.5$\,ms an under-opened
configuration suggested; the matched-$N$ verifier is therefore near-parity with
\binius{}'s, and the \zinc{} advantage is \emph{asymptotic} (the $O(N)$ vs
$\sqrt N$ split), not a moderate-$N$ headline. Second, the \zinc{} commit has a
\emph{warm} floor of $\sim$$15.5$\,ms at $\nu{=}16$ but reads $\sim$$18$\,ms in
the after-prove profile: a cold-cache penalty (the prior prove's working set
evicts the commit's encode/scatter buffers), not GPU-dispatch warmth or
reallocation. The penalty falls on repeated/throughput proving, not on a single
cold proof.
